\appendixitem{紧群的表示}

设 \(G\) 是一个紧拓扑群，且 \(\rho : G \to \mathbf{GL}(W)\) 是 \(G\) 在有限维复向量空间上的连续表示。我们称 \(\rho\) 是实型、复型或四元数型不可约表示，当且仅当对于 \(\mathbf{C}^{(G)}\)-模 \(W\)（相对于代数 \(A = \mathbf{R}^{(G)}\)），情形也是如此。设 \(H\) 是 \(W\) 上一个分离的正埃尔米特型，且在 \(G\) 下不变。

命题 3. 假设 \(\rho\) 是不可约的。

\begin{enumerate}
\def\labelenumi{\alph{enumi})}
\item
  表示 \(\rho\) 属于实型，当且仅当 W 上存在一个非零的、在 G 下不变的对称双线性型 B。在此情形，型 B 是分离的；由满足 \(w \in W\) 且 \(\mathrm{H}(w, x) = \mathrm{B}(w, x)\) 对所有 \(x \in W\) 成立的元素构成的集合 V，是 W 上一个在 G 下不变的 \(\mathbf{R}\)-结构。
\item
  表示 \(\rho\) 属于复型，当且仅当 W 上不存在在 G 下不变的非零双线性型。
\item
  表示 \(\rho\) 属于四元数型，当且仅当 W 上存在一个在 G 下不变的非零交错双线性型；这样的型必然是分离的。
\end{enumerate}

对于 \(\theta \in \mathrm{Hom}_{\mathbf{C}^{(\mathrm{G})}}(\mathrm{W},\overline{\mathrm{W}})\) 和 \(x,y\in \mathrm{W}\)，令 \(\mathrm{B}_{\theta}(x,y) = \mathrm{H}(\theta x,y)\)。则 \(\mathrm{B}_{\theta}\) 是 \(\mathrm{W}\) 上的双线性型，在 \(\mathrm{G}\) 下不变，并且当 \(\theta\) 非零时是分离的。记 \(\mathcal{B}(\mathrm{W})^{\mathrm{G}}\) 为 \(\mathrm{W}\) 上在 \(\mathrm{G}\) 下不变的双线性型空间；映射 \(\theta \mapsto \mathrm{B}_{\theta}\) 从 \(\mathrm{Hom}_{\mathbf{C}^{(\mathrm{G})}}(\mathrm{W},\overline{\mathrm{W}})\) 到 \(\mathcal{B}(\mathrm{W})^{\mathrm{G}}\) 是 \(\mathbf{C}\)-向量空间的一个同构。特别地，这推出断言 \(b\)）。

令 \(\theta\) 是一个 \(\mathbf{C}^{(\mathrm{G})}\)-同构，从 W 到 \(\overline{\mathrm{W}}\)，满足 \(\theta \circ \theta = \alpha_{\mathrm{W}}\)，其中 \(\alpha \in \{-1, +1\}\)（命题 2）；由于 \(\mathcal{B}(\mathrm{W})^{\mathrm{G}}\) 的维数为 1，存在 \(\varepsilon \in \mathbf{C}\)，使得

\[
\mathrm{B} _ {\theta} (y, x) = \varepsilon \mathrm{B} _ {\theta} (x, y) \quad \text { for   all } x, y \text { in } \mathrm{W}.
\]

反复交换变量，得到 \(\mathrm{B}_{\theta}(y,x)=\varepsilon\mathrm{B}_{\theta}(x,y)=\varepsilon^{2}\mathrm{B}_{\theta}(y,x)\)，因此 \(\varepsilon^{2}=1\)，且 \(\varepsilon\in\{-1,+1\}\)。此外，对于 W 中的 x，

\[
\mathrm{H} (\theta x, \theta x) = \mathrm{B} _ {\theta} (x, \theta x) = \varepsilon \mathrm{B} _ {\theta} (\theta x, x) = \varepsilon \mathrm{H} (\theta \circ \theta (x), x) = \varepsilon \alpha \mathrm{H} (x, x)
\]

所以由于 H 是正的，有 \(\varepsilon\alpha > 0\)，亦即 \(\varepsilon = \alpha\)。于是断言 a) 和 c) 由命题 2 得出。

记 G 上总质量为 1 的 Haar 测度为 dg。

引理 1. 设 \(W^{G}\) 是 W 的子空间，由在 G 下不变的元素组成。W 的自同态 \(\int_{G}\rho(g)dg\) 是一个投影，其像为 \(W^{G}\)，并且与 G 的作用相容。特别地，

\[
\dim \mathrm{W} ^ {\mathrm{G}} = \int_ {\mathrm{G}} \operatorname{Tr} \rho (g) d g.
\]

令 \(p = \int_{\mathrm{G}} \rho(g) dg\)；对于 \(h \in \mathrm{G}\)，

\[
\rho (h) \circ p = \int_ {\mathrm{G}} \rho (h g) d g = \int_ {\mathrm{G}} \rho (g) d g = p
\]

类似地，\(p \circ \rho(h) = p\)。因此，p 与 G 的作用相容，且其像包含于 \(W^{G}\)。若 \(w \in W^{G}\)，则 \(p(w) = \int_{\mathbf{G}} \rho(g) w \, dg = w\)，引理得证。

引理 2. 设 \(u\) 是有限维向量空间 \(E\)（定义在域 \(K\) 上）的一个自同态。则

\[
\operatorname{Tr} u ^ {2} = \operatorname{Tr} \mathbf {S} ^ {2} (u) - \operatorname{Tr} \Lambda^ {2} (u).
\]

设 \(\chi_u(\mathrm{X}) = \prod_{i=1}^{n} (\mathrm{X} - \alpha_i)\) 是在 K 的某个适当扩域中将 \(u\) 的特征多项式分解为线性因子的一个分解。我们有 \(\operatorname{Tr} u^2 = \sum_i \alpha_i^2\)、\(\operatorname{Tr} \Lambda^2(u) = \sum_{i < j} \alpha_i \alpha_j\)、\(\operatorname{Tr} \mathbf{S}^2(u) = \sum_{i \leq j} \alpha_i \alpha_j\)（参见~《代数》，第 VII 章，§5，节 5，推论 3），结果得证。

命题 4. 假设 \(\rho\) 是不可约的。那么，\(\rho\) 属于实型（分别为复型、四元数型），当且仅当积分 \(\int_{\mathrm{G}} \mathrm{Tr} \rho(g^2) dg\) 等于 1（分别为 0、-1）。

记 \(\check{\rho}\) 为 \(\rho\) 在 \(W^{*}\) 上的反步表示（定义为 \(\check{\rho}(g)={}^{t}\rho(g^{-1})\)）。将引理 2 应用于 \(\check{\rho}(g)\)，并在 G 上积分，得到

\[
\int_ {\mathrm{G}} \operatorname{Tr} \rho (g ^ {2}) d g = \int_ {\mathrm{G}} \operatorname{Tr} ^ {t} \rho (g ^ {- 2}) d g = \int_ {\mathrm{G}} \operatorname{Tr} \mathbf {S} ^ {2} (\check {\rho} (g)) d g - \int_ {\mathrm{G}} \operatorname{Tr} \Lambda^ {2} (\check {\rho} (g)) d g
\]

故由引理 1，

\[
\int_ {\mathrm{G}} \mathrm{Tr} \rho (g ^ {2}) d g = \dim (\mathbf {S} ^ {2} \mathrm{W} ^ {*}) ^ {\mathrm{G}} - \dim (\Lambda^ {2} \mathrm{W} ^ {*}) ^ {\mathrm{G}}.
\]

但是，\(S^{2}W^{*}\)（分别为 \(\Lambda^{2}W^{*}\)）可与 W 上的对称（分别为交错）双线性型空间相识别。因此，命题立即由命题 3 得出。

